\documentclass[handout]{beamer}
\usetheme{metropolis}
\usepackage[spanish,activeacute]{babel}
\usepackage{amsmath, amssymb}
\usepackage{graphicx}
\usepackage{tikz}
\usetikzlibrary{shapes.geometric, arrows.meta, positioning, fit, decorations.pathreplacing, calc, external, shapes.misc}
% \tikzexternalize[prefix=tikz/] % Descomentar para externalizar figuras TikZ y acelerar compilación
\usepackage{booktabs} % Para tablas de calidad profesional
\usepackage{multirow}
\usepackage{fontspec}
\usepackage{hyperref}

% --- Configuración de fuentes (asegúrate de tenerlas instaladas) ---
\setsansfont{Lato}
\setmonofont{Inconsolata}

% --- Configuración de Biblatex para APA ---
\usepackage[backend=biber, style=apa, sorting=nyt, dateabbrev=false]{biblatex}
\addbibresource{bio.bib} % Asegúrate de tener un archivo bio.bib en tu carpeta

% --- Metadatos de la presentación ---
\title{Retrieval Augmented Generation (RAG)}
\author{Jorge Ortiz Fuentes}
\date{\today}

% --- Configuraciones adicionales de Metropolis y Beamer ---
\metroset{block=fill} % Estilo de bloques llenos
\setbeamertemplate{navigation symbols}{} % Ocultar símbolos de navegación
\setbeamertemplate{caption}[numbered]    % Numerar captions de figuras/tablas

% --- Pie de página personalizado ---
\setbeamertemplate{footline}{
  \begin{beamercolorbox}[wd=\paperwidth,ht=2.5ex,dp=1ex]{frametitle}
    \hfill\insertframenumber/\inserttotalframenumber\hspace{3mm}
  \end{beamercolorbox}
}

% -------------------------------------------------------------------
%   METADATOS / CONFIGURACIONES
% -------------------------------------------------------------------
\setbeamercolor{normal text}{fg=black,bg=white}
\setbeamercolor{alerted text}{fg=red!80!black}
\setbeamercolor{example text}{fg=green!50!black}
\setbeamertemplate{navigation symbols}{}
\setbeamertemplate{caption}[numbered]
\setbeamertemplate{footline}{
  \begin{beamercolorbox}[wd=\paperwidth,ht=2.5ex,dp=1ex]{frametitle}
    \vspace{1pt}
    \hfill \insertframenumber/\inserttotalframenumber \hspace{3mm}
  \end{beamercolorbox}
}
\metroset{block=fill}
% -------------------------------------------------------------------

\begin{document}

% --- Diapositiva de Título ---
\begin{frame}[plain]
  \titlepage
\end{frame}

% --- Diapositiva de Contenido (Índice) ---
\begin{frame}{Contenido}
  \tableofcontents[hideallsubsections] % Oculta subsecciones para un índice más limpio en la primera vista
\end{frame}

% ===================================================================
\section{Introducción}
% ===================================================================

\begin{frame}
  \frametitle{¿Qué es RAG?}
  \begin{block}{Definición}
    \begin{itemize}
      \item RAG (\textit{Retrieval-Augmented Generation}) combina LLMs preentrenados con recuperación de información externa
      \item Los LLMs estándar se basan en conocimiento estático de su entrenamiento
      \item RAG recupera fragmentos relevantes y los usa para enriquecer la respuesta \parencite{Lewis2020RAG}
    \end{itemize}
  \end{block}
  \begin{alertblock}{Importancia}
      Mejora la factualidad, la relevancia y permite citar fuentes
  \end{alertblock}
\end{frame}

\begin{frame}
  \frametitle{Mejoras y problemas persistentes en LLMs}
  \begin{block}{Avances recientes}
    \begin{itemize}
      \item \textit{Instruction tuning} 
      \item \textit{Reinforcement learning from human feedback (RLHF)} 
      \item \textit{Ingeniería de Prompting}
    \end{itemize}
  \end{block}
  \begin{alertblock}{Desafíos actuales}
    \begin{itemize}
      \item \textit{Hallucination} (invención de información)
      \item \textit{Attribution} (rastreo del origen de la información)
      \item \textit{Staleness} (conocimiento desactualizado)
      \item \textit{Revisions} y \textit{Customization} (actualización y adaptación)
    \end{itemize}
  \end{alertblock}
  \begin{center}
    RAG surge como una solución para estos desafíos
  \end{center}
\end{frame}

\begin{frame}
  \frametitle{Paradigmas del conocimiento en LM}
  \begin{block}{Formas de conceptualizar el conocimiento}
    \begin{itemize}
      \item \textbf{\textit{Closed book} vs. \textit{Open book}} \parencite{Roberts2020Knowledge}:
            \begin{itemize}
                \item \textit{Closed book}: LM estándar memoriza conocimiento en sus parámetros
                \item \textit{Open book}: RAG accede a una base de datos externa para obtener conocimiento
            \end{itemize}
      \item \textbf{\textit{Parametric} vs. \textit{Semi-parametric}} \parencite{Yogatama2021Adaptive}:
            \begin{itemize}
                \item \textit{Parametric}: Conocimiento reside enteramente en los parámetros del modelo
                \item \textit{Semi-parametric}: RAG combina LLM (paramétrico) con un \textit{retriever} (no paramétrico)
            \end{itemize}
    \end{itemize}
  \end{block}
\end{frame}

% ===================================================================
\section{Orígenes y evolución del RAG}
% ===================================================================

\begin{frame}
  \frametitle{Formalización y antecedentes}
  \begin{block}{Hito principal}
    \begin{itemize}
      \item RAG fue formalizado y popularizado por \textcite{Lewis2020RAG}
      \item Demostraron superioridad en tareas intensivas en conocimiento
      \item Tiene raíces en sistemas de QA anteriores
    \end{itemize}
  \end{block}
\end{frame}

\begin{frame}
  \frametitle{Contexto: \textit{Open-Domain Question Answering} (ODQA)}
  \begin{block}{El desafío de ODQA}
    \begin{itemize}
      \item Responder preguntas sobre una amplia gama de temas, basándose en un gran corpus de conocimiento (ej. Wikipedia)
      \item Modelos pre-RAG enfrentaban limitaciones:
        \begin{itemize}
            \item Conocimiento estático (limitado a sus datos de entrenamiento)
            \item Tendencia a alucinar (inventar respuestas)
            \item Dificultad para actualizar conocimiento o incorporar información nueva
            \item Poca interpretabilidad sobre el origen de la respuesta
        \end{itemize}
    \end{itemize}
  \end{block}
\end{frame}

\begin{frame}
  \frametitle{Novedades del trabajo de \textcite{Lewis2020RAG}}
  \begin{block}{Arquitectura y contribuciones}
    \begin{itemize}
      \item Propusieron explícitamente la arquitectura RAG para tareas intensivas en conocimiento
      \item Combinaron un \textit{retriever} no paramétrico con un generador paramétrico (\textit{seq2seq}, BART)
      \item El \textit{retriever} busca información relevante en Wikipedia y el generador condiciona su respuesta
    \end{itemize}
  \end{block}
  \begin{alertblock}{Impacto}
    \begin{itemize}
        \item Mejora significativa en la factualidad de las respuestas
        \item Reducción de alucinaciones al basar la generación en evidencia explícita
    \end{itemize}
  \end{alertblock}
\end{frame}

% ===================================================================
\section{Arquitectura del RAG}
% ===================================================================

\begin{frame}
  \frametitle{Flujo habitual de RAG}
  \begin{figure}
    \centering
    \begin{tikzpicture}[
        scale=0.70,
        transform shape,
        font=\normalsize,
        node distance=1cm, % General node distance
        >=Stealth
      ]

      %% —— Bloque rojo (Docs → Doc Encoder → Retriever → Query Encoder) ——
      \node[
        draw, fill=red!20,
        rounded corners,
        text width=3em,
        align=center,
        minimum height=2em
      ] (docs) {Docs};

      \node[
        draw, fill=red!20,
        rounded corners,
        text width=4em,
        align=center,
        minimum height=2em,
        right=1cm of docs
      ] (doc_encoder) {Doc Encoder};

      \node[
        cylinder, shape border rotate=90,
        draw, fill=red!20,
        aspect=0.5,
        minimum height=2em,
        minimum width=1.5em,
        text centered,
        right=1cm of doc_encoder
      ] (retriever_db) {Retriever};

      \node[
        draw, fill=red!20,
        rounded corners,
        text width=4em,
        align=center,
        minimum height=2em,
        above=0.5cm of retriever_db
      ] (query_encoder) {Query Encoder};

      \draw[->] (docs)        -- (doc_encoder);
      \draw[->] (doc_encoder) -- (retriever_db);
      \draw[->] (query_encoder) -- (retriever_db);

      %% —— Bloque naranja (Input, Prompt, Context) a la derecha ——
      \node[
        draw, fill=orange!20,
        rounded corners,
        text width=3em,
        align=center,
        minimum height=2em,
        % MODIFIED POSITIONING:
        above right=0.4cm and 1.2cm of query_encoder
      ] (input_gen) {Input};

      \node[
        draw, fill=orange!20,
        rounded corners,
        text width=3em,
        align=center,
        minimum height=2em,
        below=0.3cm of input_gen
      ] (prompt_gen) {Prompt};

      \node[
        draw, fill=orange!20,
        rounded corners,
        text width=3em,
        align=center,
        minimum height=2em,
        below=0.3cm of prompt_gen
      ] (context_gen) {Context};

      %% —— Generator y Output ——
      \node[
        draw, fill=blue!20,
        rounded corners,
        text width=4em,
        align=center,
        minimum height=7.5em,
        right=1.5cm of prompt_gen
      ] (generator_gen) {Generator};

      \node[
        draw, fill=purple!20,
        rounded corners,
        text width=3em,
        align=center,
        minimum height=2em,
        right=1cm of generator_gen
      ] (output_gen) {Output};

      %% —– Flechas internas ——
      \draw[->] (input_gen)   -- (generator_gen);
      \draw[->] (prompt_gen)  -- (generator_gen);
      \draw[->] (context_gen) -- (generator_gen);
      \draw[->] (generator_gen) -- (output_gen);

      %% —— Conexiones cruzadas ——
      % De Input a QueryEncoder (codo automático con “-|”)
      \draw[->] (input_gen.west) -| (query_encoder.north);
      % De Retriever a Context
      \draw[->] (retriever_db.east) -| (context_gen.south);

    \end{tikzpicture}
    \caption{Arquitectura RAG (adaptada de clase de Stanford CS25 V3 sobre \textit{Retrieval Augmented Language Models})}
  \end{figure}
\end{frame}

\begin{frame}
  \frametitle{Componentes centrales de RAG}
  \begin{block}{Recuperador (Retriever)}
    \begin{itemize}
      \item Codificación de documentos y consultas en vectores (embeddings).
      \item Construcción de un índice \textit{Approximate Nearest Neighbor (ANN)} para permitir búsquedas eficientes.
      \item Gestión de un datastore con pares clave-valor que asocia cada embedding con su fragmento de texto o documento completo.
      \item Proceso de consulta: codificar la pregunta, realizar búsqueda ANN para obtener los \(k\) fragmentos más relevantes, y aplicar post-procesamiento (re-ranking, filtrado) si es necesario.
    \end{itemize}
  \end{block}
\end{frame}

\begin{frame}
  \frametitle{Componentes centrales de RAG}
  \begin{block}{Generador (Generator)}
  \begin{itemize}
    \item LLM preentrenado, generalmente basado en la arquitectura Transformer.
    \item Se encarga de generar la respuesta final basándose en la consulta y la información recuperada.
    \item Los generadores se pueden clasificar en dos categorías:
  \begin{itemize}
    \item \textbf{Generadores por defecto:} La mayoría de los LLMs preentrenados, como los de las series GPT, Mistral y Gemini. 
    \item \textbf{Generadores aumentados por recuperación (RA):} Modelos preentrenados que incluyen módulos específicos para fusionar la información recuperada, como RETRO \parencite{Borgeaud2021Improving}.
  \end{itemize}
  \end{itemize}
  \end{block}
\end{frame}

% ===================================================================
\section{Ventajas de RAG}
% ===================================================================

\begin{frame}
  \frametitle{Beneficios clave de RAG (I)}
  \begin{block}{Conocimiento actualizado y específico}
    \begin{itemize}
        \item Acceso a información no presente en datos de entrenamiento (ej. datos recientes, de dominio específico)
        \item No se requiere reentrenar (\textit{customization}, combate \textit{staleness})
    \end{itemize}
  \end{block}
  \begin{block}{Reducción de alucinaciones}
    \begin{itemize}
      \item Fundamentar respuestas en evidencia recuperada disminuye significativamente la invención de información y reduce alucinaciones \parencite{Lewis2020RAG, Shuster2021Retrieval, Dai2024RetrievalAugmented}
    \end{itemize}
  \end{block}
\end{frame}

\begin{frame}
  \frametitle{Beneficios clave de RAG (II)}
  \begin{block}{Transparencia y explicabilidad}
    \begin{itemize}
        \item Permite citar fuentes, facilitando la verificación y comprensión del origen de la información (\textit{attribution})
        \item Aumenta la confianza e interpretabilidad
    \end{itemize}
  \end{block}
  \begin{block}{Actualización y control}
    \begin{itemize}
        \item Actualizar la base de conocimiento externa es más económico y rápido que reentrenar un LLM (\textit{revisions})
        \item Mayor control sobre la información utilizada, facilitando personalización
    \end{itemize}
  \end{block}
\end{frame}

% ===================================================================
\section{Desventajas y desafíos de RAG}
% ===================================================================

\begin{frame}
  \frametitle{Limitaciones y obstáculos}
  \begin{alertblock}{Dependencia de la calidad de recuperación}
    \begin{itemize}
        \item La efectividad de RAG está ligada a la calidad del \textit{retriever}
        \item Información irrelevante o incorrecta en la recuperación afecta la respuesta final \parencite{Dai2024RetrievalAugmented}
    \end{itemize}
  \end{alertblock}
  \begin{block}{Latencia y complejidad}
    \begin{itemize}
        \item El proceso de recuperación introduce latencia adicional
        \item Implementar y mantener pipelines RAG robustos puede ser complejo
    \end{itemize}
  \end{block}
\end{frame}

\begin{frame}
  \frametitle{Otros desafíos importantes}
  \begin{block}{Sesgos y costos}
    \begin{itemize}
        \item Posibilidad de heredar y amplificar sesgos de la base de conocimiento externa
        \item Costos computacionales significativos (recuperación e inferencia con prompts grandes)
    \end{itemize}
  \end{block}
   \begin{alertblock}{Manejo de información}
    \begin{itemize}
        \item Dificultad para manejar fragmentos recuperados que son irrelevantes, contradictorios o ruidosos
    \end{itemize}
  \end{alertblock}
\end{frame}

% ===================================================================
\section{Componentes clave y técnicas avanzadas}
% ===================================================================
\begin{frame}
  \frametitle{Profundizando en los mecanismos de RAG}
    \begin{block}{Áreas de optimización e investigación}
        La efectividad de RAG depende de la sofisticación de sus componentes y las técnicas empleadas.
        \begin{itemize}
            \item El recuperador (\textit{retriever})
            \item El generador (\textit{generator})
            \item La ingeniería de \textit{prompts}
        \end{itemize}
    \end{block}
\end{frame}

% -------------------------------------------------------------------
\subsection{Recuperador (\textit{Retriever})}
% -------------------------------------------------------------------
\begin{frame}
  \frametitle{\textit{Retriever}: tipos de \textit{embeddings} y técnicas (I)}
  \begin{columns}[T]
    \begin{column}{0.5\textwidth}
      \begin{block}{\textit{Sparse Retrieval}}
        \begin{itemize}
            \item Basados en conteo de palabras
            \item Bueno para coincidencias exactas
            \item Ejemplos: TF-IDF, BM25
          \end{itemize}
      \end{block}
    \end{column}
    \begin{column}{0.5\textwidth}
      \begin{block}{\textit{Dense Retrieval}}
        \begin{itemize}
            \item Capturan similitud semántica
            \item Puede fallar en palabras clave exactas
            \item Ejemplos: Sentence Transformers, OpenAI Embeddings
        \end{itemize}
      \end{block}
    \end{column}
  \end{columns}
  \begin{block}{\textit{Hybrid Search}}
    Suele ser la mejor opción práctica, combinando fortalezas de \textit{sparse} y \textit{dense retrieval}
  \end{block}
\end{frame}

\begin{frame}
  \frametitle{\textit{Retriever}: bases de datos vectoriales}
  \begin{block}{Bases de datos vectoriales: almacenamiento y búsqueda}
      \begin{itemize}
        \item Una base de datos vectorial está diseñada para almacenar, gestionar y buscar datos codificados como vectores.
        \item Se componen de tres elementos clave:
            \begin{itemize}
                \item \textbf{Chunking y Encoding:} Los documentos se dividen en fragmentos (\textit{chunks}) y se codifican como vectores.
                \item \textbf{Indexing:} Se construye un índice especializado para acelerar la búsqueda, utilizando algoritmos de ANN (por ejemplo, IVFPQ, HNSW, Annoy).
                \item \textbf{Datastore (base de datos vectorial):} Almacena los datos originales (ej. el texto de los \textit{chunks}) asociados a sus vectores, a menudo como pares clave-valor.
            \end{itemize}
      \end{itemize}
  \end{block}
\end{frame}

\begin{frame}
  \frametitle{\textit{Retriever}: Bases de datos vectoriales (II)}
  \begin{block}{Bases de datos vectoriales open source}
    \begin{itemize}
      \item \textbf{FAISS}:  
      Biblioteca de Meta en C++/Python para búsquedas de similitud en vectores a gran escala, con soporte GPU.

      \item \textbf{Annoy}:  
      Biblioteca basada en árboles de partición aleatorios (random projection trees) en C++/Python, ligeros y mmap.

      \item \textbf{Qdrant}:  
      Base de datos con motor en Rust que ofrece búsqueda de similitud, cuantización y persistencia de vectores para entornos escalables.

      \item \textbf{pgvector}:  
      Extensión de PostgreSQL que añade un tipo vectorial y funciones ANN para RAG.
    \end{itemize}
  \end{block}
\end{frame}

\begin{frame}
  \frametitle{\textit{Retriever}: Re-ranking para mejorar la precisión}
  \begin{block}{¿Qué es el re-ranking?}
    Es una segunda etapa de filtrado que refina los resultados del Retrieval para asegurar que se seleccionen los documentos más relevantes.
  \end{block}
  \begin{block}{Proceso en dos fases}
    \begin{enumerate}
        \item \textbf{Recuperación inicial (\textit{Retriever})}:
        Un modelo rápido recupera un conjunto amplio de documentos potencialmente relevantes (ej. los 20-100 mejores). Se prioriza la \textbf{exhaustividad (\textit{recall})}.
        
        \item \textbf{Re-clasificación (\textit{Re-ranker})}:
        Un modelo más potente y lento evalúa la relevancia de cada documento recuperado con respecto a la consulta y los reordena. Se prioriza la \textbf{precisión} \parencite{humeau2020polyencoderstransformerarchitecturespretraining}.
    \end{enumerate}
  \end{block}
\end{frame}

\begin{frame}
  \frametitle{Ejemplo de re-ranking}
  \begin{figure}[h!]
    \centering
    \includegraphics[scale=0.45]{pics/rerank_cohere.png}
    \caption{Ejemplo de re-ranking en RAG (extraído de post de Cohere sobre Rerank)}
  \end{figure}
\end{frame}

% -------------------------------------------------------
\subsection{Generador (\textit{Generator})}
% -------------------------------------------------------------------
\begin{frame}
  \frametitle{Generador: síntesis de Información}
    \begin{block}{Rol del generador}
        \begin{itemize}
            \item Típicamente un LLM potente (GPT, Claude, Gemini, DeepSeek, etc.)
            \item Sintetiza una respuesta coherente y precisa basada en la consulta original y la información recuperada
            \item Capacidad de integrar y razonar sobre la información recuperada es crucial
        \end{itemize}
    \end{block}
\end{frame}

\begin{frame}
  \frametitle{Ejemplo de prompt para RAG}
  \begin{block}{Prompt ejemplo}
    \footnotesize\ttfamily
    You are an assistant for question-answering tasks. Use the following retrieved context to answer the question.\\
    If you don’t know the answer, just say that you don’t know. Keep it under three sentences.\\[0.5em]
    \textbf{Question:} \{question\}\\
    \textbf{Context:} \{context\}\\
    \textbf{Answer:}
    \normalfont
  \end{block}
  \vspace{0.5em}
  \raggedleft
  \small{Fuente: \href{https://smith.langchain.com/hub/rlm/rag-prompt}{LangChain Prompt Hub}}
\end{frame}


% -------------------------------------------------------------------
\subsection{Ingeniería de \textit{prompts} para RAG}
% -------------------------------------------------------------------
\begin{frame}
  \frametitle{Elaboración de \textit{Prompts} Efectivos}
    \begin{block}{Impacto del \textit{prompt}}
        \begin{itemize}
            \item La estructura del \textit{prompt} (consulta + chunks recuperados) afecta significativamente la calidad de la respuesta
            \item Un buen \textit{prompt} guía al LLM a enfocarse, citar fuentes, manejar contradicciones \parencite{Zhang2024Towards}
        \end{itemize}
    \end{block}
\end{frame}

% ===================================================================
\section{Panorama actual: discusiones y debates}
% ===================================================================
\begin{frame}
    \frametitle{Estado del arte}
    El campo de RAG está en evolución y desarrollo.
\end{frame}

% -------------------------------------------------------------------
\subsection{RAG vs. Ventanas de Contexto Largas (LCWs)}
% -------------------------------------------------------------------

\begin{frame}
  \frametitle{RAG vs. Ventanas de Contexto Largas (LCWs)}
\begin{figure}[h!]
	\centering
	\includegraphics[scale=0.15]{pics/gemini_lcw_meme.png}
\end{figure}
\end{frame}

\begin{frame}
  \frametitle{RAG vs. \textit{Long Context Windows} (LCWs)}
  \begin{columns}[T]
    \begin{column}{0.5\textwidth}
      \begin{block}{LCWs}
        \begin{itemize}
            \item Conceptualmente simples si el conocimiento cabe en el \textit{prompt}
            \item Pueden sufrir de \textit{lost in the middle} \parencite{liu2023lostmiddlelanguagemodels}
            \item Más costoso
            \item No resuelven conocimiento dinámico o muy masivo
        \end{itemize}
      \end{block}
    \end{column}
    \begin{column}{0.5\textwidth}
      \begin{block}{RAG}
        \begin{itemize}
            \item Ventajoso para bases de conocimiento masivas y dinámicas
            \item Mayor precisión en algunos casos
            \item Control granular y eficiencia de costos
        \end{itemize}
      \end{block}
    \end{column}
  \end{columns}
  \begin{center}
    RAG y LCWs no son mutuamente excluyentes ya que se pueden combinar.
  \end{center}
\end{frame}

\begin{frame}
  \frametitle{El problema de ``Lost in the Middle''}
  \begin{columns}
    \column{0.4\textwidth}
    \begin{block}{Investigación de \textcite{liu2023lostmiddlelanguagemodels}}
      \begin{itemize}
          \item LLMs tienden a prestar más atención a la información al principio y al final del contexto largo
          \item La información en el medio a menudo se ignora o se pierde
      \end{itemize}
    \end{block}

    \column{0.6\textwidth}
    \begin{figure}
        \centering
        \begin{tikzpicture}[scale=0.8, transform shape]
            \draw[->] (0,0) -- (6,0) node[right] {Posición};
            \draw[->] (0,0) -- (0,4) node[above] {Precisión};
            \draw[smooth, thick, blue] plot coordinates {(0.5,3.5) (1.5,3) (3,1.5) (4.5,2.8) (5.5,3.3)};
            \node at (3, -0.5) {Inicio \hspace{2cm} Medio \hspace{2cm} Fin};
            \node[text width=4.5cm, align=center] at (3,4.5) {\small Curva de rendimiento en forma de U};
        \end{tikzpicture}
    \end{figure}
  \end{columns}
\end{frame}

% -------------------------------------------------------------------
\subsection{Tipos de RAG}
% -------------------------------------------------------------------
\begin{frame}
  \frametitle{Taxonomía de Sistemas RAG}
  \begin{block}{Niveles de implementación (basado en \textcite{Dai2024RetrievalAugmented})}
    \begin{itemize}
        \item \textbf{\textit{Naive RAG}}:
            \begin{itemize}
                \item Implementación básica: indexación, recuperación, generación
                \item \textit{Chunking} simple, \textit{embeddings} genéricos
            \end{itemize}
        \item \textbf{\textit{Advanced RAG}}:
            \begin{itemize}
                \item Optimizaciones en pre/post-procesamiento y recuperación (ej. re-ranking)
            \end{itemize}
        \item \textbf{\textit{RAG modular y agéntico}}:
            \begin{itemize}
                \item Arquitecturas flexibles, módulos intercambiables (ej. Self-RAG \parencite{Asai2023SelfRAG})
            \end{itemize}
    \end{itemize}
  \end{block}
\end{frame}


% ===================================================================
\section{Evaluación de sistemas RAG}
% ===================================================================
\begin{frame}
    \frametitle{Midiendo el rendimiento de RAG}
    El RAG puede evaluarse tanto desde el \textit{retrieval} como desde la \textit{generación}.
\end{frame}

% -------------------------------------------------------------------
\subsection{Métricas para el Retrieval}
% -------------------------------------------------------------------
\begin{frame}
  \frametitle{Recall@K}
  \begin{block}{Definición}
    \begin{itemize}
        \item Mide la proporción de documentos relevantes que son recuperados entre los K primeros resultados
        \item Es importante cuando se quiere asegurar que la mayoría de los documentos relevantes sean encontrados
    \end{itemize}
  \end{block}
  \begin{block}{Fórmula}
    \[ \text{Recall@K} = \frac{\text{N° de docs recuperados en el top K}}{\text{Número total de documentos relevantes}} \]
  \end{block}
\end{frame}

\begin{frame}
  \frametitle{Precision@K}
  \begin{block}{Definición}
    \begin{itemize}
        \item Mide la proporción de documentos recuperados en los K primeros resultados que son realmente relevantes
        \item Es crucial cuando la exactitud de los primeros resultados es prioritaria
    \end{itemize}
  \end{block}
  \begin{block}{Fórmula}
    \[ \text{Precision@K} = \frac{\text{N° de docs recuperados en el top K}}{\text{K}} \]
  \end{block}
\end{frame}

\begin{frame}
  \frametitle{Hit Rate@K}
  \begin{block}{Definición}
    \begin{itemize}
        \item Indica la proporción de consultas para las cuales al menos un documento relevante se encuentra entre los K primeros resultados recuperados
        \item Es una métrica útil para entender si el sistema suele encontrar algo relevante rápidamente
    \end{itemize}
  \end{block}
  \begin{block}{Fórmula}
    \[ \text{Hit Rate@K} = \frac{1}{|Q|} \sum_{i=1}^{|Q|} \mathbb{I}(\text{hit}_i@K) \]
    Donde \(|Q|\) es el número total de consultas y \(\mathbb{I}(\text{hit}_i@K)\) es 1 si la consulta \(i\) tiene al menos un resultado relevante en el top K, y 0 en caso contrario.
  \end{block}
\end{frame}

\begin{frame}
  \frametitle{Mean Reciprocal Rank (MRR)}
  \begin{block}{Definición}
    \begin{itemize}
        \item Evalúa la calidad de un sistema de recuperación basándose en la posición del primer resultado relevante
        \item Es el promedio de la inversa del rango del primer resultado correcto para un conjunto de consultas
    \end{itemize}
  \end{block}
  \begin{block}{Fórmula}
    \[ \text{MRR} = \frac{1}{|Q|} \sum_{i=1}^{|Q|} \frac{1}{\text{rank}_i} \]
    Donde \(|Q|\) es el número total de consultas y \(\text{rank}_i\) es la posición del primer documento relevante para la consulta \(i\). Si no hay documentos relevantes para una consulta, \(\frac{1}{\text{rank}_i}\) se considera 0 para esa consulta.
  \end{block}
\end{frame}

% -------------------------------------------------------------------
\subsection{Métricas para el Generator}
% -------------------------------------------------------------------
\begin{frame}
  \frametitle{Evaluación del componente de generación}
    \begin{block}{Dimensiones de calidad}
        \begin{itemize}
            \item \textbf{Fidelidad/Veracidad (\textit{Faithfulness/Factuality})}: adhesión al contexto recuperado
            \item \textbf{Relevancia de la respuesta (\textit{Answer Relevance})}: pertinencia a la consulta del usuario
            \item \textbf{Coherencia y fluidez}: calidad lingüística de la respuesta
            \item \textbf{No toxicidad / sesgos}: ausencia de lenguaje ofensivo o sesgos
        \end{itemize}
    \end{block}
\end{frame}


% ===================================================================
\section{RAG Multimodal (MM-RAG)}
% ===================================================================
\begin{frame}
  \frametitle{Extendiendo RAG más allá del texto}
    \begin{block}{Concepto de MM-RAG}
        \begin{itemize}
            \item El conocimiento humano es multimodal (texto, imágenes, audio, video)
            \item MM-RAG busca extender el paradigma RAG para recuperar y razonar sobre información de múltiples modalidades
        \end{itemize}
    \end{block}
    \begin{alertblock}{Desafíos significativos}
        \begin{itemize}
            \item Creación de \textit{embeddings} multimodales efectivos
            \item Diseño de arquitecturas de recuperación para datos modales múltiples
            \item Capacidad de LLMs para fusionar y generar desde información multimodal
        \end{itemize}
    \end{alertblock}
\end{frame}

\begin{frame}
\frametitle{Arquitectura de RAG Multimodal}
El flujo de MM-RAG adapta el pipeline tradicional para manejar datos heterogéneos.

\begin{enumerate}
    \item \textbf{Consulta multimodal}: el usuario puede preguntar usando texto, una imagen, o una combinación de ambas.
    \item \textbf{Recuperación multimodal}:
    \begin{itemize}
        \item Se usan \textbf{embeddings multimodales} para codificar la consulta y los documentos (imágenes, texto, audio) en un espacio semántico compartido.
        \item Se buscan los fragmentos más relevantes en una base de datos vectorial multimodal.
    \end{itemize}
    \item \textbf{Generación multimodal}: un \textbf{LLM Multimodal (MLLM)} procesa la consulta y el contexto recuperado para generar una respuesta.
\end{enumerate}
\end{frame}

\begin{frame}
\frametitle{Componente clave: Embeddings Multimodales}
\begin{block}{Definición}
Son representaciones vectoriales que capturan la semántica de diferentes modalidades en un espacio compartido \parencite{radford2021learningtransferablevisualmodels}.
\end{block}
\begin{columns}[T]
\begin{column}{0.5\textwidth}
\begin{block}{¿Cómo se crean?}
\begin{itemize}
    \item Generalmente con aprendizaje contrastivo a gran escala.
    \item Se entrena un modelo para maximizar la similitud entre pares correctos (ej. imagen y su leyenda) y minimizarla para pares incorrectos.
\end{itemize}
\end{block}
\end{column}
\begin{column}{0.5\textwidth}
\begin{block}{Modelos notables}
\begin{itemize}
    \item \textbf{CLIP} (OpenAI): Pionero en alinear texto e imágenes.
    \item \textbf{ImageBind} (Meta): Unifica 6 modalidades (texto, imagen/video, audio, etc.) en un solo espacio.
\end{itemize}
\end{block}
\end{column}
\end{columns}
\end{frame}

\begin{frame}
\frametitle{Componente clave: LLMs Multimodales (MLLM)}
\begin{block}{¿Qué son los MLLMs?}
Son LLMs extendidos para procesar, comprender y generar contenido a partir de múltiples modalidades.
\end{block}
\begin{block}{Arquitectura y modelos líderes}
\begin{itemize}
    \item \textbf{Arquitectura}: usualmente, un codificador de modalidad (ej. un Vision Transformer) proyecta la información no textual a un espacio que el LLM puede entender.
    \item \textbf{Empresas y modelos}:
    \begin{itemize}
        \item \textbf{OpenAI}: GPT-4V, GPT-4o y modelos en adelante
        \item \textbf{Google}: Gemini 1.5 en adelante
        \item \textbf{Meta}: Llama 3.2 (con capacidades multimodales)
    \end{itemize}
\end{itemize}
\end{block}
\end{frame}


% ===================================================================
\section{Frameworks y herramientas}
% ===================================================================
\begin{frame}
  \frametitle{Ecosistema para el desarrollo de RAG}
    \begin{block}{Herramientas populares}
        \begin{itemize}
            \item \textbf{LangChain}: framework más popular para construir aplicaciones LLM, incluyendo pipelines RAG complejos
            \item \textbf{LlamaIndex}: framework especializado la preparación e indexación de datos
            \item \textbf{Haystack}: framework end‐to‐end listo para producción de RAG/búsqueda semántica con énfasis en escalabilidad y monitorización.
        \end{itemize}
    \end{block}
\end{frame}

% ===================================================================
\section{Direcciones futuras e investigación abierta}
% ===================================================================
\begin{frame}
  \frametitle{Algunas lineas de investigación actuales (I)}
        \begin{itemize}
            \item \textbf{RAG Agéntico}: Combinar RAG con agentes LLM autónomos para problemas complejos (ej: Deep Research)
            \item \textbf{GraphRAG}: Usar grafos de conocimiento como fuente externa para razonar sobre relaciones
            \item \textbf{Auto-mejora y aprendizaje por retroalimentación}: Sistemas RAG que aprenden de errores (ej. \textcite{Asai2023SelfRAG})
        \end{itemize}
\end{frame}

\begin{frame}
  \frametitle{Algunas lineas de investigación actuales (II)}
        \begin{itemize}
            \item Optimización de costos, latencia y escalabilidad para bases de conocimiento enormes
            \item Uso de RAG para diferentes task de NLP
            \item Mejoras en robustez y fiabilidad: manejo de información ruidosa, contradictoria o desactualizada
        \end{itemize}
\end{frame}


% --- Diapositiva de Referencias ---
\begin{frame}[allowframebreaks] 
  \frametitle{Referencias}
  \printbibliography
\end{frame}

\end{document}


% TO DO: agregar prompt ejemplo de RAG y explicar métricas con más detalles
